%% Chapter 6 — Delta Analysis
\chapter{Delta Analysis}
\label{ch:delta}

\section{Introduction}

This chapter presents the core output of the \BDE: the parcel-by-parcel and aggregate biodiversity unit delta between the 2016 baseline (\Tzero) and the 2026 present-day assessment (\Tone). All calculations follow the methodology established in Chapter~\ref{ch:methodology}, with full uncertainty quantification via Monte Carlo simulation.


\section{Aggregate Delta Summary}

\begin{table}[H]
\centering
\caption{Aggregate biodiversity unit delta summary.}
\label{tab:aggregate-delta}
\begin{tabular}{lrrr}
\toprule
\textbf{Metric} & \textbf{\Tzero{} (2016)} & \textbf{\Tone{} (2026)} & \textbf{Delta} \\
\midrule
Total area-based BU                          & \dataplaceholder{T0} & \dataplaceholder{T1} & \dataplaceholder{DELTA} \\
95\% CI lower bound                           & --- & --- & \dataplaceholder{CI LOW} \\
95\% CI upper bound                           & --- & --- & \dataplaceholder{CI HIGH} \\
Net percentage change                         & --- & --- & \dataplaceholder{PCT}\% \\
Parcels with positive delta                   & --- & --- & \dataplaceholder{N POS} \\
Parcels with negative delta                   & --- & --- & \dataplaceholder{N NEG} \\
Parcels with zero delta ($|\Delta| < 0.01$)   & --- & --- & \dataplaceholder{N ZERO} \\
\bottomrule
\end{tabular}
\end{table}

\begin{confidencebox}{HIGH}
The aggregate delta of \dataplaceholder{DELTA BU} (\dataplaceholder{PCT}\%) is statistically significant at the 95\% confidence level, as the confidence interval [\dataplaceholder{CI LOW}, \dataplaceholder{CI HIGH}] does not span zero.
\end{confidencebox}


\section{Parcel-Level Delta Distribution}

\begin{figure}[H]
\centering
\begin{tikzpicture}
  \draw[dreamlab-light,fill=dreamlab-light,rounded corners] (0,0) rectangle (12,7);
  \node[dreamlab-grey,font=\sffamily\large] at (6,3.5) {Figure placeholder: Delta histogram};
  \node[dreamlab-grey,font=\sffamily\footnotesize] at (6,2.5) {../output/figures/delta\_histogram.png};
\end{tikzpicture}
\caption{Distribution of per-parcel $\Delta$BU values. Positive values indicate biodiversity gain; negative values indicate loss. The vertical dashed line marks zero.}
\label{fig:delta-histogram}
\end{figure}

\begin{table}[H]
\centering
\caption{Descriptive statistics of the per-parcel $\Delta$BU distribution.}
\label{tab:delta-stats}
\begin{tabular}{lr}
\toprule
\textbf{Statistic} & \textbf{Value (BU)} \\
\midrule
Mean $\Delta$BU per parcel        & \dataplaceholder{MEAN} \\
Median $\Delta$BU per parcel      & \dataplaceholder{MEDIAN} \\
Standard deviation                & \dataplaceholder{SD} \\
Minimum (largest loss)            & \dataplaceholder{MIN} \\
Maximum (largest gain)            & \dataplaceholder{MAX} \\
Interquartile range               & \dataplaceholder{IQR} \\
Skewness                          & \dataplaceholder{SKEW} \\
\bottomrule
\end{tabular}
\end{table}


\section{Spatial Pattern of Change}

\begin{figure}[H]
\centering
\begin{tikzpicture}
  \draw[dreamlab-light,fill=dreamlab-light,rounded corners] (0,0) rectangle (12,8);
  \node[dreamlab-grey,font=\sffamily\large] at (6,4) {Figure placeholder: Delta map (diverging colour scale)};
  \node[dreamlab-grey,font=\sffamily\footnotesize] at (6,3) {../output/figures/delta\_map\_diverging.png};
\end{tikzpicture}
\caption{Parcel-level $\Delta$BU map using a diverging colour scale: green = gain, white = zero, red = loss. Hatched parcels indicate classification change between epochs.}
\label{fig:delta-map}
\end{figure}


\section{Delta Decomposition by Driver}

The total delta can be decomposed into contributions from three BM\,4.0 components: classification change (affecting distinctiveness), condition change, and strategic significance change:

\begin{equation}
\Delta\text{BU}_i = \underbrace{\Delta D_i \cdot \bar{C}_i \cdot \bar{S}_i \cdot A_i}_{\text{Classification effect}} + \underbrace{\bar{D}_i \cdot \Delta C_i \cdot \bar{S}_i \cdot A_i}_{\text{Condition effect}} + \underbrace{\bar{D}_i \cdot \bar{C}_i \cdot \Delta S_i \cdot A_i}_{\text{Strategic effect}} + \epsilon_i
\label{eq:delta-decomposition}
\end{equation}

\noindent where $\bar{X}$ denotes the mean of $X$ across the two epochs and $\epsilon_i$ captures interaction terms.

\begin{table}[H]
\centering
\caption{Delta decomposition by driver.}
\label{tab:delta-decomposition}
\begin{tabular}{lrr}
\toprule
\textbf{Driver} & \textbf{$\Delta$BU Contribution} & \textbf{\% of Total Delta} \\
\midrule
Habitat classification change   & \dataplaceholder{DELTA} & \dataplaceholder{PCT}\% \\
Condition change                & \dataplaceholder{DELTA} & \dataplaceholder{PCT}\% \\
Strategic significance change   & \dataplaceholder{DELTA} & \dataplaceholder{PCT}\% \\
Interaction terms               & \dataplaceholder{DELTA} & \dataplaceholder{PCT}\% \\
\midrule
\textbf{Total}                  & \dataplaceholder{TOTAL} & 100.0\% \\
\bottomrule
\end{tabular}
\end{table}

\begin{confidencebox}{MEDIUM}
The decomposition assumes that driver effects are approximately additive. In parcels where both classification and condition have changed substantially, the interaction term $\epsilon_i$ may be non-trivial. Such parcels are flagged in the per-parcel data tables (Appendix~\ref{app:data}).
\end{confidencebox}


\section{BU Calculation Worked Examples}

\subsection{Example 1: Grassland Improvement}

Parcel TOID \dataplaceholder{TOID}, area \dataplaceholder{A}\,ha:

\begin{align}
\text{BU}_{T_0} &= A \times D_{T_0} \times C_{T_0} \times S_{T_0} \notag \\
  &= \dataplaceholder{A} \times 2 \times 1.5 \times 1.00 = \dataplaceholder{BU0} \label{eq:example1-t0} \\[8pt]
\text{BU}_{T_1} &= A \times D_{T_1} \times C_{T_1} \times S_{T_1} \notag \\
  &= \dataplaceholder{A} \times 4 \times 2.5 \times 1.10 = \dataplaceholder{BU1} \label{eq:example1-t1} \\[8pt]
\Delta\text{BU} &= \dataplaceholder{BU1} - \dataplaceholder{BU0} = +\dataplaceholder{DELTA} \notag
\end{align}

This parcel transitioned from modified grassland (g4, Low distinctiveness, Fairly Poor condition) to other neutral grassland (g3c, Medium distinctiveness, Fairly Good condition), with a strategic significance upgrade from Low to Medium due to the revised LNRS draft layer.

\subsection{Example 2: Woodland Condition Decline}

Parcel TOID \dataplaceholder{TOID}, area \dataplaceholder{A}\,ha:

\begin{align}
\text{BU}_{T_0} &= \dataplaceholder{A} \times 6 \times 3.0 \times 1.15 = \dataplaceholder{BU0} \\[8pt]
\text{BU}_{T_1} &= \dataplaceholder{A} \times 6 \times 2.0 \times 1.15 = \dataplaceholder{BU1} \\[8pt]
\Delta\text{BU} &= \dataplaceholder{BU1} - \dataplaceholder{BU0} = -\dataplaceholder{DELTA}
\end{align}

Habitat classification unchanged (lowland mixed deciduous woodland, w1f), but condition declined from Good to Moderate, reflecting reduced NDVI heterogeneity consistent with canopy homogenisation or understorey loss.

\subsection{Example 3: Development Loss}

Parcel TOID \dataplaceholder{TOID}, area \dataplaceholder{A}\,ha:

\begin{align}
\text{BU}_{T_0} &= \dataplaceholder{A} \times 4 \times 2.0 \times 1.00 = \dataplaceholder{BU0} \\[8pt]
\text{BU}_{T_1} &= \dataplaceholder{A} \times 0 \times 0 \times 1.00 = 0.00 \\[8pt]
\Delta\text{BU} &= 0.00 - \dataplaceholder{BU0} = -\dataplaceholder{DELTA}
\end{align}

Parcel converted from other neutral grassland (g3c) to developed land (u1b), representing a complete loss of biodiversity value.


\section{Sensitivity Analysis}

\subsection{One-At-A-Time (OAT) Analysis}

Each input parameter was varied across its plausible range while holding all others at their nominal values:

\begin{table}[H]
\centering
\caption{Sensitivity analysis: parameter ranges and effect on total $\Delta$BU.}
\label{tab:sensitivity}
\begin{tabular}{p{40mm}ccc}
\toprule
\textbf{Parameter} & \textbf{Low} & \textbf{Nominal} & \textbf{High} \\
\midrule
Condition threshold shift ($\pm$1 SD) & \dataplaceholder{LOW} & \dataplaceholder{NOM} & \dataplaceholder{HIGH} \\
Classification probability cutoff      & \dataplaceholder{LOW} & \dataplaceholder{NOM} & \dataplaceholder{HIGH} \\
Area measurement error ($\pm$2\%)      & \dataplaceholder{LOW} & \dataplaceholder{NOM} & \dataplaceholder{HIGH} \\
LNRS layer version (v0.2 vs v0.3)     & \dataplaceholder{LOW} & \dataplaceholder{NOM} & \dataplaceholder{HIGH} \\
\bottomrule
\end{tabular}
\end{table}

\begin{figure}[H]
\centering
\begin{tikzpicture}
  \draw[dreamlab-light,fill=dreamlab-light,rounded corners] (0,0) rectangle (12,7);
  \node[dreamlab-grey,font=\sffamily\large] at (6,3.5) {Figure placeholder: Tornado diagram};
  \node[dreamlab-grey,font=\sffamily\footnotesize] at (6,2.5) {../output/figures/tornado\_sensitivity.png};
\end{tikzpicture}
\caption{Tornado diagram showing the relative influence of each parameter on the total $\Delta$BU. Bars indicate the range of total delta when each parameter is varied from low to high.}
\label{fig:tornado}
\end{figure}


\section{Monte Carlo Uncertainty Quantification}

\subsection{Simulation Design}

A Monte Carlo simulation with $N = 10{,}000$ iterations was performed, perturbing the three uncertainty sources described in Section~\ref{sec:bm4scoring} simultaneously at each iteration. The simulation produces a probability distribution of the total $\Delta\text{BU}_{\text{total}}$.

\subsection{Results}

\begin{figure}[H]
\centering
\begin{tikzpicture}
  \draw[dreamlab-light,fill=dreamlab-light,rounded corners] (0,0) rectangle (12,7);
  \node[dreamlab-grey,font=\sffamily\large] at (6,3.5) {Figure placeholder: Monte Carlo delta distribution};
  \node[dreamlab-grey,font=\sffamily\footnotesize] at (6,2.5) {../output/figures/monte\_carlo\_delta\_distribution.png};
\end{tikzpicture}
\caption{Monte Carlo simulation results: probability density of the total $\Delta$BU across 10,000 iterations. The vertical red lines mark the 2.5th and 97.5th percentiles (95\% CI); the vertical blue line marks the nominal estimate.}
\label{fig:mc-distribution}
\end{figure}

\begin{table}[H]
\centering
\caption{Monte Carlo simulation summary statistics.}
\label{tab:mc-stats}
\begin{tabular}{lr}
\toprule
\textbf{Statistic} & \textbf{Value (BU)} \\
\midrule
Nominal $\Delta$BU          & \dataplaceholder{NOM} \\
Monte Carlo mean             & \dataplaceholder{MEAN} \\
Monte Carlo std.~dev.        & \dataplaceholder{SD} \\
2.5th percentile (95\% CI)  & \dataplaceholder{P025} \\
97.5th percentile (95\% CI) & \dataplaceholder{P975} \\
Probability of net gain      & \dataplaceholder{P GAIN}\% \\
Probability of net loss      & \dataplaceholder{P LOSS}\% \\
\bottomrule
\end{tabular}
\end{table}

\subsection{Convergence Diagnostics}

\begin{figure}[H]
\centering
\begin{tikzpicture}
  \draw[dreamlab-light,fill=dreamlab-light,rounded corners] (0,0) rectangle (12,5);
  \node[dreamlab-grey,font=\sffamily\large] at (6,2.5) {Figure placeholder: MC convergence trace};
  \node[dreamlab-grey,font=\sffamily\footnotesize] at (6,1.5) {../output/figures/mc\_convergence.png};
\end{tikzpicture}
\caption{Monte Carlo convergence trace: running mean and 95\% CI width as a function of iteration count. Convergence (CI width change $<$0.1\% per 100 iterations) is achieved by iteration \dataplaceholder{N CONVERGE}.}
\label{fig:mc-convergence}
\end{figure}

\begin{confidencebox}{HIGH}
The Monte Carlo simulation converged within \dataplaceholder{N CONVERGE} iterations. The 95\% CI width of \dataplaceholder{CI WIDTH} BU is \dataplaceholder{PCT}\% of the nominal delta, indicating acceptable precision for statutory reporting purposes.
\end{confidencebox}


\section{Comparison with 10\% BNG Threshold}

While this analysis is not a formal BNG assessment for a planning application, it is instructive to compare the observed delta against the statutory 10\% net gain threshold:

\begin{equation}
\text{10\% BNG threshold} = 0.10 \times \text{BU}_{T_0} = 0.10 \times \dataplaceholder{T0 BU} = \dataplaceholder{THRESHOLD} \text{ BU}
\label{eq:bng-threshold}
\end{equation}

The observed delta of \dataplaceholder{DELTA BU} is \dataplaceholder{ABOVE/BELOW} this threshold. If the study area were subject to a BNG obligation, the current trajectory would \dataplaceholder{MEET/NOT MEET} the 10\% requirement \dataplaceholder{WITH/WITHOUT} additional habitat creation or enhancement measures.

\begin{confidencebox}{SPECULATIVE}
This comparison is illustrative only. A formal BNG assessment would require consideration of habitat trading rules, the mitigation hierarchy, temporal discounting for habitat creation lag, and the specific development footprint, none of which are modelled here.
\end{confidencebox}
